% 第4章 4.6：本章对象关系图。局部可积函数 -> 广义函数 -> 广义导数；有局部可积表示时为弱导数；加上 L^p 控制得 W^{k,p}；
% 下方四个问题（完备性、平滑稠密、嵌入与紧性、迹与零边界）及其所需条件。旁注三个检验例子。
\begin{tikzpicture}[font=\figurefont\footnotesize,
  box/.style={draw=FigureBlue,fill=FigureBlueFill,rounded corners=1.8pt,line width=.7pt,align=center,inner sep=3pt,minimum height=8mm,minimum width=22mm},
  q/.style={draw=FigureGreen,fill=white,rounded corners=1.8pt,line width=.7pt,align=center,inner sep=2.5pt,minimum width=24mm,minimum height=11mm,font=\figurefont\scriptsize},
  ex/.style={font=\FigureNoteFont\scriptsize,text=FigureMuted,align=left},
  arr/.style={-{Latex[length=1.8mm]},draw=FigureStroke,line width=.7pt}]
\node[box] (loc) at (0,0) {局部可积函数\\$f\in L^1_{\mathrm{loc}}$};
\node[box] (gen) at (4.0,0) {广义函数\\$T_f$};
\node[box] (der) at (8.0,0) {广义导数\\$D^\alpha T_f$};
\node[box] (weak) at (8.0,-1.65) {弱导数\\$D^\alpha T_f=T_g$};
\node[box] (sob) at (4.0,-1.65) {Sobolev空间\\$W^{k,p}$};
\draw[arr] (loc) -- node[above,font=\figurefont\scriptsize]{测试配对} (gen);
\draw[arr] (gen) -- node[above,font=\figurefont\scriptsize]{总可定义} (der);
\draw[arr] (der) -- node[right,font=\figurefont\scriptsize,align=left]{有局部\\可积表示} (weak);
\draw[arr] (weak) -- node[above,font=\figurefont\scriptsize]{$L^p$控制} (sob);
\node[q] (q1) at (-0.4,-3.5) {完备性\\Banach/Hilbert};
\node[q] (q2) at (2.6,-3.5) {平滑稠密\\整域或到边界};
\node[q] (q3) at (5.6,-3.5) {嵌入与紧性\\有界Lipschitz域};
\node[q] (q4) at (8.6,-3.5) {迹、零边界\\与Poincaré};
\foreach \n in {q1,q2,q3,q4} \draw[arr] (sob.south) -- (\n.north);
\node[ex,anchor=west] at (9.6,0.0) {$D\delta_0$：非测度};
\node[ex,anchor=west] at (9.6,-1.3) {$|x|$：有，$\operatorname{sgn}x$};
\node[ex,anchor=west] at (9.6,-1.75) {$H$：无，得$\delta_0$};
\node[ex,anchor=west] at (9.6,-2.2) {Cantor：无，奇异测度};
\end{tikzpicture}
